\label{anlage:sekundaerbefunde}

Die Studie \textit{ReUse in der Bauindustrie stärken} von BaselCircular und Cirkla (2025) basiert auf 18 explorativen Expert:inneninterviews und einer anschließenden Umfrage unter 94 Akteur:innen der Schweizer Bauwirtschaft. Die nachfolgend zusammengefassten Ergebnisse werden als ergänzende Sekundärbefunde herangezogen; sie stammen nicht aus einer eigenen Erhebung.

\subsection{Relevante Befunde}

\subsubsection{Frühe Berücksichtigung von ReUse}

Die Studie nennt die frühe Einbindung von ReUse in den Planungsprozess als wichtige Voraussetzung für eine breitere Anwendung. Dazu zählen eine projektbezogene ReUse-Strategie, die rechtzeitige Erfassung verfügbarer Bauteile und die Identifikation geeigneter Quellobjekte. Wiederverwendbare Komponenten müssen bekannt sein, bevor grundlegende Planungsentscheidungen getroffen werden.

\subsubsection{Verfügbarkeit von Bauteilinformationen}

Als Hemmnis wird die geringe Verfügbarkeit belastbarer Angaben zu Quellobjekten und Bauteilen beschrieben. Genannt werden unter anderem Verbindungen, Lebensdauer, Verfügbarkeit, Tragfähigkeit und statische Eigenschaften. Fehlende oder unvollständige Informationen erschweren die Beurteilung von Ausbau, Prüfung und erneuter Verwendung.

\subsubsection{Digitale Inventare und standardisierte Daten}

Digitale Bauteilinventare und Potenzialanalysen werden als Grundlage für Rückbau, Vermittlung und Wiedereinbau eingeordnet. Ergänzend verweist die Studie auf Materialpässe, Gebäudemodelle, Rückbauinventare und den Schweizer Inventarstandard Swiss-Inv. Einheitliche Datenstrukturen sollen den Informationsaustausch zwischen Inventarisierung, Planung, Ausschreibung und Dokumentation erleichtern.

\subsubsection{Ökologische Bewertung in frühen Planungsphasen}

Die Studie empfiehlt, Wiederverwendung bereits in frühen Ökobilanzen zu berücksichtigen und Inventardaten mit Bewertungs- und Ausschreibungssystemen zu verknüpfen. Ein konkretes Berechnungsverfahren für einzelne gebrauchte Bauteile wird nicht beschrieben.

\subsection{Quelle}

Ackermann, Sarah; Kliem, Daniel; Kellenberger, Daniel; Olender, Margarete; Held, Samuel (2025): \textit{ReUse in der Bauindustrie stärken. Studie zur aktuellen Entwicklung der Wiederverwendung in der Bauindustrie und zum Potenzial ihrer Skalierung}. BaselCircular; Cirkla, Basel.
